\documentclass[11pt,a4paper]{article}
\usepackage[margin=2.5cm]{geometry}
\usepackage{setspace}
\usepackage{parskip}
\usepackage{enumitem}
\usepackage{cite}
\setstretch{1.1}
\begin{document}
\begin{center}
    {\LARGE \textbf{Master Thesis Proposal}}\\
    {\large \textbf{PF Generation: Can an LLM pipeline with simulated-student-driven refinement generate Productive Failure exercises that approach expert quality?}}
\end{center}
\vspace{0.4cm}
\noindent\textit{Historical proposal, retained for provenance. The phases and evaluations below are proposed work, not completed experiments. The current thesis chapters describe the method and evidence through October 10, 2026. Prior PF generation by Rahimi et al.~\cite{rahimi2025} qualifies the original novelty claim.}

\section*{Background}
Productive Failure (PF) is a well-established learning design in which students attempt challenging, ill-structured problems before receiving instruction, generating diverse representations and solution methods (RSMs) that prepare them for deeper learning \cite{kapur2008,sinha2021}. However, designing effective PF exercises remains labor-intensive, requiring careful calibration of ill-structuredness across multiple dimensions \cite{kapur2009}. Recent work has explored using large language models (LLMs) to automatically generate educational content, particularly multiple-choice questions and distractors \cite{alhazmi2024}. Anticipating student approaches remains distinct from generating mathematically correct content \cite{feng2024,parikh2025}. Rahimi et al.~\cite{rahimi2025} already study automatic PF problem generation. The open question here is whether a pipeline can elicit meaningful approach diversity and support productive failure, not whether PF generation has ever been attempted.
\section*{Purpose of this Paper}

We propose an LLM-based pipeline that automatically generates PF exercises 
and evaluates their quality through simulated student interactions. 
The proposed research program has three phases:

\begin{description}
    \setlength{\itemsep}{0.5em}
    
    \item[\textbf{(1) Exercise generation}] where an LLM transforms well-structured problems into ill-structured ones guided by Kapur's  five criteria of ill-structuredness \cite{kapur2009}, including parameter complexity, solution complexity, problem space constraints accessibility, and collaboration potential, using retrieval-augmented few-shot examples \cite{puech2025}.
    \item[\textbf{(2) Simulated student evaluation}] where student profiles defined by Knowledge Components (KCs) from the Knowledge-Learning-Instruction framework \cite{lu2024} attempt the generated exercises at varying ability levels, yielding measures of RSM diversity and productive failure rates. To ensure the diversity signal reflects genuine response variation rather than superficial rewording, we induce diversity through three complementary mechanisms: (i) KC-based persona variation \cite{lu2024}, (ii) grounding in real student response distributions from large-scale math question banks \cite{wang2021eedi}, and (iii) stochastic sampling across runs.
    \item[\textbf{(3) Iterative refinement}] where simulated student signals feed back into the generation process to improve exercise quality \cite{madaan2023selfrefine}.
\end{description}
\section*{Data, Evaluation and Baselines}

\paragraph{Proposed data sources.} The proposal identifies expert-authored PF exercises collected through StratL \cite{puech2025} and the Eedi platform \cite{wang2021eedi}. The scale of the external response dataset is not the number of responses used in this repository's experiments; current subsets and their distinct roles are reported in the thesis.

\paragraph{Proposed evaluation.} The original plan would evaluate final exercises on held-out seed problems along two tiers:
\begin{description}
    \setlength{\itemsep}{0.3em}
    \item[\textbf{(i) Automatic metrics}] including mathematical validity via a symbolic solver, an independent LLM-as-judge \cite{zheng2023judging} on Kapur's five criteria using a different model and rubric than the in-loop judge of Phase~3 to detect optimization leakage, an embedding-based diversity score over a held-out sentence encoder \cite{tevet2021diversity}, and a memorization check via n-gram overlap against the retrieval pool \cite{carlini2023memorization}.
    \item[\textbf{(ii) Human expert evaluation}] by learning scientists and teachers using Likert ratings and pairwise preference comparisons.
\end{description}

\paragraph{Proposed baselines.} The original plan compares zero-shot, few-shot, and retrieval-augmented generation on a common seed-to-PF interface, with expert-authored StratL exercises as design references rather than a measured upper bound. Proposed ablations remove refinement or KC-based personas. These student-refinement ablations have not been completed. Later generation comparisons instead reconstruct the prompt setting of Rahimi et al.~\cite{rahimi2025} and compare complete direct and decomposed methods.
 
\bibliographystyle{unsrt}
\bibliography{references}
 
\end{document}